% % sec 4.3.6

% \subsection{Quaternions e justificativa de sua utilização}
% \label{sec:representacao_quaternion}
 
% Os parâmetros de Euler, ou quatérnions, constituem uma alternativa para representar a atitude da espaçonave, pois evitam as singularidades cinemáticas associadas aos ângulos de Euler e permitem descrever a orientação de forma regular ao longo da integração numérica \cite{wie2008space}.
% Um quaternion unitário é definido como um elemento do espaço quaternional de quatro dimensões, pertencente ao conjunto
% dos quaternions unitários $q = [w,\vec{v}\,]$, com $\|q\| = 1$, em que $w$ é a parte escalar e $\vec{v}$ é a parte vetorial. Cada quaternion unitário $q$ corresponde a uma rotação de ângulo $\theta$ em torno de um eixo unitário $\vec{u}$, dada por:

% \begin{equation}
% q(\theta,\vec{u}) = \cos\!\left(\frac{\theta}{2}\right) + \sin\!\left(\frac{\theta}{2}\right)\,\vec{u}.
% \label{eq_quat_axis_angle}
% \end{equation}

% Essa parametrização fornece uma representação livre de singularidades geométricas para a integração contínua das equações diferenciais de atitude.
% Em aplicações numéricas, impõe-se a condição de unidade $\|q\|=1$ e utiliza-se o desvio de norma como critério de monitoramento de erros de integração ao longo da simulação.

% Do ponto de vista físico, os quaternions preservam a norma e realizam rotações por multiplicação, assegurando uma transição suave entre orientações sucessivas.
% Embora exijam uma variável adicional em comparação aos ângulos de Euler, os quatérnions apresentam como principal vantagem a ausência de singularidades cinemáticas. Além disso, sua forma diferencial é particularmente adequada à integração numérica da atitude \cite{wie2008space}.

% %%%%%%%%%%%%%%%%%%%%%%%%%
% \subsubsection{Definição e componentes do quaternion}
% \label{sec:componentes_quaternion}

% O espaço quaternional é um espaço vetorial de quatro dimensões, composto por um elemento escalar e três elementos imaginários mutuamente ortogonais $\{i,j,k\}$, que obedecem às regras de multiplicação introduzidas por Hamilton (1844):
% \begin{equation}
% i^2 = j^2 = k^2 = ijk = -1,
% \label{eq_hamilton_quadrados}
% \end{equation}
% \begin{equation}
% ij = k,\quad jk = i,\quad ki = j,\quad
% ji = -k,\quad kj = -i,\quad ik = -j.
% \label{eq_hamilton_produtos}
% \end{equation}

% Um quaternion qualquer pode ser escrito como:
% \begin{equation}
% q = w + x\,i + y\,j + z\,k.
% \label{eq_quaternion_forma}
% \end{equation}

% Ou, de forma vetorial compacta:

% \begin{equation}
% q = [\,w,\vec{v}\,] = [\,w,x,y,z\,].
% \label{eq_quaternion_wxyz}
% \end{equation}

% Em que $w$ é a parte escalar e $\vec{v} = [x,y,z]$ é a parte vetorial.
% Neste trabalho será adotada a notação $q = [w,x,y,z]^T$.
% A magnitude (ou norma) de um quaternion é definida por:

% \begin{equation}
% \|q\| = \sqrt{w^2 + x^2 + y^2 + z^2}.
% \label{eq_norma_quaternion}
% \end{equation}

% E o quaternion unitário é obtido normalizando-se $q$ de modo que $\|q\| = 1$.
% Em aplicações de atitude, a condição de unidade constitui uma restrição do estado e pode ser empregada para monitorar o acúmulo de erros numéricos durante a integração.
% Desvios de $\|q\|$ em relação a $1$ indicam erro numérico e podem ser corrigidos por renormalização do quaternion, evitando que a representação de atitude se degrade ao longo do tempo.

% A relação entre os quaternions e a velocidade angular do corpo é obtida a partir da equação cinemática diferencial \cite{wie2008space}.
% Considerando o vetor de estado quaternion:

% \begin{equation}
% \mathbf{q} =
% \begin{bmatrix}
% w \\
% x \\
% y \\
% z
% \end{bmatrix},
% \qquad
% \vec{\omega} =
% \begin{bmatrix}
% \omega_x \\ \omega_y \\ \omega_z
% \end{bmatrix}.
% \label{eq_q_omega_vetores}
% \end{equation}

% A equação cinemática pode ser escrita na forma matricial:

% \begin{equation}
% \begin{bmatrix}
% \dot{w} \\[2pt]
% \dot{x} \\[2pt]
% \dot{y} \\[2pt]
% \dot{z}
% \end{bmatrix}
% =
% \frac{1}{2}
% \begin{bmatrix}
% 0 & -\omega_x & -\omega_y & -\omega_z \\
% \omega_x & 0 & \omega_z & -\omega_y \\
% \omega_y & -\omega_z & 0 & \omega_x \\
% \omega_z & \omega_y & -\omega_x & 0
% \end{bmatrix}
% \begin{bmatrix}
% w \\ x \\ y \\ z
% \end{bmatrix}.
% \label{eq_cinematica_quaternion_wxyz}
% \end{equation}

% De forma compacta, pode-se escrever:

% \begin{equation}
% \dot{\mathbf{q}} = \tfrac{1}{2}\,\boldsymbol{\Omega}(\vec{\omega})\,\mathbf{q}.
% \label{eq_quat_compacta}
% \end{equation}

% Em que $\boldsymbol{\Omega}(\vec{\omega})$ é a matriz $4\times 4$ definida em \eqref{eq_cinematica_quaternion_wxyz}.

%%%%%%%%%%% após 20/04/2026


% sec 4.3.5

\subsection{Quaternions e justificativa de sua utilização}
\label{sec:representacao_quaternion}

Os parâmetros de Euler, ou quatérnions, constituem uma alternativa para representar a atitude da espaçonave, pois evitam as singularidades cinemáticas associadas aos ângulos de Euler e permitem descrever rotações finitas de forma regular ao longo da integração numérica \cite{wie2008space}.
Diferentemente dos ângulos de Euler, os quatérnions não são empregados neste trabalho como coordenadas generalizadas independentes, mas como parâmetros de atitude sujeitos à restrição de norma unitária. Assim, sua utilização não altera a formulação física da dinâmica e da cinemática orbital anteriormente estabelecida, mas apenas fornece uma parametrização mais conveniente para a propagação da orientação da espaçonave.
Um quaternion unitário pode ser escrito como:

\begin{equation}
\mathbf{q}=
\begin{bmatrix}
q_0 \\
q_1 \\
q_2 \\
q_3
\end{bmatrix}.
\qquad
q_0^2+q_1^2+q_2^2+q_3^2=1,
\label{eq_quaternion_unitario}
\end{equation}

Em que $q_0$ é a parte escalar e $\begin{bmatrix} q_1 & q_2 & q_3 \end{bmatrix}^{T}$ corresponde à parte vetorial. A restrição algébrica em \eqref{eq_quaternion_unitario} mostra que os quatro elementos do quaternion não são independentes.
Geometricamente, um quaternion unitário representa uma rotação de ângulo $\theta$ em torno de um eixo unitário $\vec{u}$, sendo expresso por:

\begin{equation}
\mathbf{q}(\theta,\vec{u})=
\begin{bmatrix}
\cos\left(\dfrac{\theta}{2}\right) \\
u_x\sin\left(\dfrac{\theta}{2}\right) \\
u_y\sin\left(\dfrac{\theta}{2}\right) \\
u_z\sin\left(\dfrac{\theta}{2}\right)
\end{bmatrix}.
\label{eq_quaternion_axis_angle}
\end{equation}

Essa parametrização é livre de singularidades geométricas e, por isso, é particularmente útil na integração contínua das equações diferenciais da atitude. Em aplicações numéricas, a condição de unidade em \eqref{eq_quaternion_unitario} é utilizada também como critério de monitoramento do erro de integração.
A matriz de cossenos diretores associada ao quaternion unitário, que relaciona o referencial orbital $A$ ao referencial do corpo $B$, pode ser escrita como \cite{wie2008space}:

\begin{equation}
\mathbf{C}^{B/A}=
\begin{bmatrix}
q_0^2+q_1^2-q_2^2-q_3^2 & 2(q_1q_2+q_0q_3) & 2(q_1q_3-q_0q_2) \\
2(q_1q_2-q_0q_3) & q_0^2-q_1^2+q_2^2-q_3^2 & 2(q_2q_3+q_0q_1) \\
2(q_1q_3+q_0q_2) & 2(q_2q_3-q_0q_1) & q_0^2-q_1^2-q_2^2+q_3^2
\end{bmatrix}.
\label{eq_dcm_quaternion}
\end{equation}

A expressão \eqref{eq_dcm_quaternion} estabelece a ligação direta entre os quatérnions e os cossenos diretores utilizados anteriormente na formulação cinemática e na expressão do torque de gradiente de gravidade. Em particular, os termos $C_{13}$, $C_{23}$ e $C_{33}$ podem ser obtidos diretamente da terceira coluna de \eqref{eq_dcm_quaternion}.
A relação entre o quaternion e a velocidade angular total da espaçonave, expressa no referencial do corpo, é dada pela equação cinemática diferencial \cite{wie2008space}:

\begin{equation}
\dot{\mathbf{q}}
=
\frac{1}{2}
\begin{bmatrix}
0 & -\omega_1 & -\omega_2 & -\omega_3 \\
\omega_1 & 0 & \omega_3 & -\omega_2 \\
\omega_2 & -\omega_3 & 0 & \omega_1 \\
\omega_3 & \omega_2 & -\omega_1 & 0
\end{bmatrix}
\mathbf{q}.
\label{eq_cinematica_quaternion_wxyz}
\end{equation}

Em que:

\begin{equation}
    \label{eq_wbn}
\vec{\omega}^{B/N}=
\begin{bmatrix}
\omega_1 \\
\omega_2 \\
\omega_3
\end{bmatrix}.
\end{equation}

Onde a Equação~\ref{eq_wbn} é a velocidade angular total da espaçonave em relação ao referencial inercial, expressa no referencial do corpo.

De forma compacta, a Equação~\eqref{eq_cinematica_quaternion_wxyz} pode ser escrita como:

\begin{equation}
\dot{\mathbf{q}}=\frac{1}{2}\,\boldsymbol{\Omega}\big(\vec{\omega}^{B/N}\big)\,\mathbf{q}.
\label{eq_quat_compacta}
\end{equation}

Em que $\boldsymbol{\Omega}\big(\vec{\omega}^{B/N}\big)$ é a matriz cinemática associada à velocidade angular total da espaçonave.
Portanto, os quatérnions fornecem uma parametrização globalmente regular da atitude, mantendo consistência com a formulação cinemática orbital adotada neste trabalho. Além disso, por meio da relação \eqref{eq_dcm_quaternion}, essa parametrização permite reescrever diretamente os cossenos diretores presentes nas equações diferenciais da cinemática e na expressão do torque de gradiente de gravidade, assegurando coerência entre as diferentes formas de representação da atitude utilizadas ao longo deste capítulo.